\section{Q1: Why does an MLP perform poorly on CIFAR-10?}
\label{sec:q1}

\subsection{(a) Missing spatial inductive bias vs.\ capacity vs.\ optimisation}

\paragraph{Hypotheses.} \textbf{\Hone\ (chosen reason): no spatial inductive bias.} The MLP's first layer gives every pixel
an independent weight to every hidden unit. It has no notion that neighbouring pixels are related (locality), and a
feature learned at one position does not transfer to another (no weight sharing). With 45k images it should therefore
\emph{overfit} with position-specific templates rather than underfit. \textbf{\Htwo: insufficient capacity.}
\textbf{\Hthree: optimisation difficulty.}

\paragraph{Design.} The core is a 2$\times$2 factorial with architecture \{MLP-M, CNN\} at matched parameters crossed with
input \{original, fixed pixel permutation\}. The permutation is one shuffle of the 1024 pixel positions (RGB values move
together), applied identically to training, validation and test images. It keeps the information in each image, the
labels, the parameter count, the FLOPs and the optimiser unchanged, and removes only spatial structure. An MLP is
indifferent to it: the shuffle is absorbed by permuting the columns of its first weight matrix, and i.i.d.\ initialisation
is permutation-invariant. A width sweep (MLP-S/M/L/XL, 34$\times$ parameters, fixed depth) tests \Htwo. Eval-mode
accuracy and loss on a fixed 5k training subset, logged every epoch, test \Htwo/\Hthree.

\paragraph{Fairness.} A comparison is fair here if the models differ only in architecture and have the same opportunity
to learn: (i) the same data, split and normalisation, with no augmentation; (ii) parameters matched within 0.2\%;
(iii) the same budget of 30 epochs at batch 128 (10,560 AdamW steps), with the same cosine schedule, weight decay 0.05, loss,
BatchNorm and ReLU; (iv) equal tuning effort, with the LR chosen from \{3e-4, 1e-3, 3e-3\} and the early-stopping epoch
chosen on validation for \emph{every} model and condition, including the permuted ones; (v) 3 seeds per model.
Compute per image is \emph{not} matched: weight sharing gives the CNN $\approx$118$\times$ more MACs
(Tab.~\ref{tab:models}). The permutation control rules this confound out, because CNN-perm has exactly the CNN's
parameters \emph{and} FLOPs.

\begin{table}[t]
\centering\footnotesize\setlength{\tabcolsep}{2.5pt}
\caption{Q1(a) predictions, written before running.}
\label{tab:q1pred}
\begin{tabular}{@{}lccc@{}}
\toprule
Measurement & \Hone\ bias & \Htwo\ capacity & \Hthree\ optimisation \\
\midrule
MLP train accuracy & high, big gap & low & low, loss stalls \\
Widening MLP 34$\times$ & small gain & large gain & --- \\
CNN at equal params & much better & similar & maybe better \\
Permutation: CNN & large drop & no change & no change \\
Permutation: MLP & no change & no change & no change \\
\bottomrule
\end{tabular}
\end{table}

\begin{table}[t]
\centering\footnotesize\setlength{\tabcolsep}{2.5pt}
\caption{Q1(a) results (mean$\pm$sd over 3 seeds). LR: selected on validation. Train: eval-mode accuracy on a fixed 5k
training subset at the final epoch. Val/Test: best-validation checkpoint. Accuracies in \%.}
\label{tab:q1res}
\begin{tabular}{@{}lrlrrr@{}}
\toprule
Model & Params & LR & Train & Val & \textbf{Test} \\
\midrule
MLP-S      & \res{q1.MLP-S.params}      & \res{q1.MLP-S.lr}      & \res{q1.MLP-S.trainfit}      & \res{q1.MLP-S.val}      & \res{q1.MLP-S.test} \\
MLP-M      & \res{q1.MLP-M.params}      & \res{q1.MLP-M.lr}      & \res{q1.MLP-M.trainfit}      & \res{q1.MLP-M.val}      & \res{q1.MLP-M.test} \\
MLP-L      & \res{q1.MLP-L.params}      & \res{q1.MLP-L.lr}      & \res{q1.MLP-L.trainfit}      & \res{q1.MLP-L.val}      & \res{q1.MLP-L.test} \\
MLP-XL     & \res{q1.MLP-XL.params}     & \res{q1.MLP-XL.lr}     & \res{q1.MLP-XL.trainfit}     & \res{q1.MLP-XL.val}     & \res{q1.MLP-XL.test} \\
CNN        & \res{q1.CNN.params}        & \res{q1.CNN.lr}        & \res{q1.CNN.trainfit}        & \res{q1.CNN.val}        & \res{q1.CNN.test} \\
\midrule
MLP-M-perm & \res{q1.MLP-M-perm.params} & \res{q1.MLP-M-perm.lr} & \res{q1.MLP-M-perm.trainfit} & \res{q1.MLP-M-perm.val} & \res{q1.MLP-M-perm.test} \\
CNN-perm   & \res{q1.CNN-perm.params}   & \res{q1.CNN-perm.lr}   & \res{q1.CNN-perm.trainfit}   & \res{q1.CNN-perm.val}   & \res{q1.CNN-perm.test} \\
\bottomrule
\end{tabular}
\end{table}

\begin{table}[t]
\centering\footnotesize\setlength{\tabcolsep}{3pt}
\caption{Q1(a) hypothesis tests: paired differences in accuracy (percentage points), on validation and on test.
Pooled seed-to-seed sd of validation accuracy: \res{q1.ev.seedsd}\,pp.}
\label{tab:q1ev}
\begin{tabular}{@{}lrr@{}}
\toprule
Quantity & Val & Test \\
\midrule
CNN $-$ MLP-M (equal parameters) & \res{q1.ev.cnnminusmlp.val} & \res{q1.ev.cnnminusmlp.test} \\
MLP-XL $-$ MLP-S (34$\times$ parameters) & \res{q1.ev.capgain.val} & \res{q1.ev.capgain.test} \\
Permutation effect on MLP-M & \res{q1.ev.permmlp.val} & \res{q1.ev.permmlp.test} \\
Permutation effect on CNN & \res{q1.ev.permcnn.val} & \res{q1.ev.permcnn.test} \\
CNN-perm $-$ MLP-M-perm & \res{q1.ev.permgap.val} & \res{q1.ev.permgap.test} \\
Interaction (CNN $-$ MLP perm effect) & \res{q1.ev.interaction.val} & \res{q1.ev.interaction.test} \\
\bottomrule
\end{tabular}
\end{table}

\begin{figure*}[t]
\centering
\resfig{fig_a1_curves_factorial.png}{0.88\textwidth}{2.45in}
\caption{Q1(a) training (eval mode, 5k subset) and validation accuracy and loss for the matched comparison and the
permutation controls (mean$\pm$sd over 3 seeds; dotted line: mean best-validation epoch).}
\label{fig:q1curves}
\end{figure*}

\begin{figure*}[t]
\centering
\resfig{fig_a3_capacity_and_factorial_report.png}{0.8\textwidth}{2.1in}
\caption{Left: validation accuracy of the MLP width sweep against parameter count, with the CNN and CNN-perm at MLP-M's
size. Right: the 2$\times$2 permutation factorial (error bars: $\pm$1 sd over 3 seeds). Redrawn from
\texttt{table\_a3\_validation\_summary.csv} by \texttt{scripts/make\_figures.py}; training/validation curves of the
capacity sweep are in \texttt{results/q1\_outputs/fig\_a2\_curves\_capacity.png}.}
\label{fig:q1cap}
\end{figure*}

\paragraph{Results.} Tables~\ref{tab:q1res} and \ref{tab:q1ev} and Figs.~\ref{fig:q1curves} and \ref{fig:q1cap} give the
measurements. The interpretation below was generated from them by the pre-registered rules in the Q1 notebook (\S6):

\restext{q1a}

\subsection{(b) Is structure learned before memorisation?}

\paragraph{Setting.} 40\% of the 45k training labels are replaced by a class drawn uniformly from the nine \emph{wrong}
classes. The draw is seeded and fixed for the whole run; validation labels stay clean. A corrupted label is unrelated to its
image, so fitting it cannot generalise. The model is \textbf{MLP-XL} (14.7M parameters, $\approx$330 per training
example; reused from (a)), trained with AdamW at a \emph{constant} LR of 1e-3, \emph{no} weight decay and 150 epochs. Weight
decay and early stopping counteract memorisation, and an LR schedule would add phase changes of its own to the trajectory. A clean-label run with the
same initialisation and batch order serves as the control~\cite{zhang2017understanding,arpit2017closer}.

\paragraph{Measurements} (whole training set in eval mode, 4 times per epoch for the first 10 epochs, then every epoch):
clean-train accuracy; \emph{memorisation rate}, the fraction of corrupted examples predicted as their assigned wrong label;
the \emph{shadow rate}, the fraction predicted as a second wrong label drawn the same way but never shown to the model. Without
example-specific memorisation the two labels are exchangeable, so the \emph{memorisation excess} (rate $-$ shadow) gets a paired $z$-score. We also record
the true-class accuracy on corrupted examples (chance 10\%), validation accuracy, and per example the epoch from which it
stays fitted.

\paragraph{Decision rule (pre-registered).} Onset $t_\text{mem}$ is the first evaluation point after which the excess stays
significant ($z>3$). We conclude that ``structure comes before memorisation'' if, at $t_\text{mem}$, validation accuracy has
covered $\geq$50\% of its rise from chance to its peak \emph{and} true-class accuracy on corrupted examples exceeds 20\%.

\begin{figure*}[t]
\centering
\resfig{fig_b1_memorisation.png}{\textwidth}{2.3in}
\caption{Q1(b): MLP-XL with 40\% corrupted labels. (a) Accuracy measures; the dashed line is the clean-label control's
validation accuracy. (b) Losses on the clean and corrupted subsets. (c) Fraction of examples stably fitted by each epoch.
Vertical lines: peak validation accuracy (grey) and memorisation onset (red). The $x$-axis is linear below 1 epoch and logarithmic above.}
\label{fig:q1b}
\end{figure*}

\begin{table}[t]
\centering\footnotesize\setlength{\tabcolsep}{2.2pt}
\caption{Q1(b) trajectory (\%). $\dagger$: peak validation accuracy; $\ast$: memorisation onset. ``True'': true class
predicted on corrupted examples; ``Mem.'': assigned wrong label predicted; ``Shad.'': never-shown wrong label predicted.}
\label{tab:q1b}
\resinput{tab_q1b_milestones}{\begin{tabular}{@{}rrrrrrr@{}}
\toprule
Epoch & Val & Val (clean run) & Clean train & True & Mem. & Shad. \\
\midrule
\multicolumn{7}{c}{\pending\ (rows generated from \texttt{table\_b1\_*.csv})}\\
\bottomrule
\end{tabular}}
\end{table}

\begin{table}[t]
\centering\footnotesize\setlength{\tabcolsep}{3pt}
\caption{Q1(b) test accuracy before and after memorisation, and the median epoch at which examples become stably fitted.}
\label{tab:q1btest}
\begin{tabular}{@{}llrrr@{}}
\toprule
Run & Checkpoint & Epoch & Val & \textbf{Test} \\
\midrule
40\% noise & early-stopped (val) & \res{q1b.test.noisy.early.ep} & \res{q1b.test.noisy.early.val} & \res{q1b.test.noisy.early.test} \\
40\% noise & final & \res{q1b.test.noisy.final.ep} & \res{q1b.test.noisy.final.val} & \res{q1b.test.noisy.final.test} \\
clean & early-stopped (val) & \res{q1b.test.clean.early.ep} & \res{q1b.test.clean.early.val} & \res{q1b.test.clean.early.test} \\
clean & final & \res{q1b.test.clean.final.ep} & \res{q1b.test.clean.final.val} & \res{q1b.test.clean.final.test} \\
\midrule
\multicolumn{5}{@{}l}{Median stably-fitted epoch: clean ex.\ \res{q1b.med.clean}, corrupted ex.\ \res{q1b.med.corrupted},} \\
\multicolumn{5}{@{}l}{same images with true labels (control run) \res{q1b.med.same}.} \\
\bottomrule
\end{tabular}
\end{table}

\paragraph{Results.} Fig.~\ref{fig:q1b} and Tabs.~\ref{tab:q1b} and \ref{tab:q1btest}. Peak validation accuracy is
\res{q1b.valpeak}\% at epoch \res{q1b.tpeak}, and memorisation onset is at epoch \res{q1b.tmem}. Pre-registered decision:
\textbf{\res{q1b.decision}}. Generated interpretation:

\restext{q1b}
